\subsection{Results: What Online Adaptation Is Worth}
\label{sec:online:results}

We evaluate on synthetic churn traces built from the same degradation model as
\S\ref{sec:weights}: hosts move between healthy, degraded and failed
regimes under a Markov process, degraded hosts draw a kind and a severity from the
five failure modes of \S\ref{sec:weights}, and within-regime drift keeps the
instantaneous argmin moving even when no regime changes. Policies are scored
against an offline optimum computed by dynamic programming over the time--host
trellis, which is exact when the primaries decouple and a feasible---hence
upper-bounding---schedule when congestion couples them. The DP is verified against
exhaustive enumeration on small instances.

Alongside the policies of \S\ref{sec:online:model} we evaluate
\textsc{failover}: the deployed static assignment, but migrating when the current
backup \emph{hard-fails}. This is the fair version of what a production system
does. It does not react to a host that is merely degraded, which is the regime the
online policies target, so comparing against it rather than against pure
\textsc{static} separates the value of reacting to death from the value of
reacting to degradation, and avoids crediting the online policies for beating a
strawman.

\paragraph{Snapshot-chasing fails exactly as predicted.}
Table~\ref{tab:online-migration} sweeps the migration cost $M$ over a $24\times$
range. \textsc{snapshot}'s excess ratio rises from $2.44$ to $28.46$, an $11.7\times$
degradation, and the reason is its migration spend:
climbs from $46\%$ to $95\%$ of its total cost. At the largest $M$ it is paying
nineteen-twentieths of its bill for the privilege of never being more than one
round out of date. This is Theorem~\ref{thm:snapshot} in the feature-based model
rather than in the adversarial construction, and it is worth stating plainly what
it implies for the rest of this paper: the Hungarian \textsc{optimal} of
\S\ref{sec:optimal} is a legitimate upper bound on \emph{one-shot} placement
quality and is not an upper bound on anything once state moves. Reported as an
achievable target for a deployed selector, it would be badly misleading.

The hysteresis policies are flat over the same sweep---$2.35$ to $2.03$ for
$\alpha=2$, drifting slightly \emph{downward}---because their threshold is
expressed in units of $M$, so a more expensive migration automatically demands
proportionally more evidence before it is taken. No retuning is required when the
replicated state grows or the link slows down, which is the practical payoff of
Proposition~\ref{prop:budget}.

\paragraph{Almost all of the available gain comes from failover, not from
degradation-chasing.}
This is the main empirical result, and it is a negative one.
Table~\ref{tab:online-churn} sweeps the churn rate with hard failures enabled.
At the most volatile setting, \textsc{static} reaches an excess ratio of $42.25$:
a policy that never moves eventually sits on dead replicas and its cost is
dominated by them. Adding failover alone drops this to $4.64$, closing $91.2\%$ of
the distance to the offline optimum. Adding full degradation-aware hysteresis on
top of failover takes it to $4.59$ --- a further $0.13$ percentage points.

Table~\ref{tab:online-persistence} isolates the effect by disabling hard failures
entirely, so that \textsc{failover} and \textsc{static} coincide and the only
question is whether to react to degradation. Across a $150\times$ range of
degradation durations the two sit between $1.08$ and $1.29$, and they trade places:
hysteresis is ahead at very short and very long durations and behind in the middle.
There is no regime in this sweep where reacting to degradation is clearly worth
its migration cost. The reason is visible in the arithmetic of the cost model. A
degraded host costs roughly $0.5$ per round more than a healthy one, and a
migration costs $M$; but degradation is transient, so a policy that leaves also
has to come back, making the round trip $2M$ against a saving of $0.5$ per round.
Migration pays only once degradation outlasts $2M/\mathrm{gap}$ rounds, which at
the $M = 3$ used here is about $12$ --- the mean degradation duration the model
was parameterised with. The policies tie because this sweep sits on its own
crossover.

That makes the value of $M$ decisive rather than incidental, and $M$ has so far
been assumed rather than measured. \S\ref{sec:online:M} measures it, finds it
several times larger, and reports what that does to the comparison: the tie does
not survive.

That is a useful thing for an operator to be able to compute. It also sharpens the
paper's wider finding rather than contradicting it: adaptive machinery again fails
to beat a simple deterministic rule, and this time there is a closed-form condition
saying when it could.

\paragraph{A weak migration budget is a real failure mode.}
Scaling the topology from 8 to 64 hosts leaves the ordering unchanged. The
well-provisioned policies are stable---\textsc{failover} improves from $2.26$ to
$1.77$ and \textsc{hyst}($2$) stays between $2.03$ and $2.30$---but
\textsc{hyst}($0.5$), whose Proposition~\ref{prop:budget} guarantee permits
migration to reach twice its service cost, degrades from $2.39$ to $34.84$. With
more candidate backups there are more ways for a marginal host to look briefly
attractive, and a policy with a loose budget takes all of them. This is the
practical content of Proposition~\ref{prop:lower}: the achievable ratio degrades
with $k$, and the degradation shows up first in the policy with the weakest
budget.

Table~\ref{tab:online-alpha} reports the budget directly. The bound holds at every
$\alpha$, as it must, but with large slack---at $\alpha = 1$ the policy is permitted
to spend as much on migration as on service and in fact spends $13.4\%$ of it. The
guarantee is therefore not what limits the policy in practice; it is what rules out
the catastrophic tail that Theorem~\ref{thm:snapshot} exhibits. Ratio is best
around $\alpha = 4$ ($1.536$) and is insensitive between $1$ and $8$, so the
parameter does not need careful tuning---in contrast to the RL penalty of
Table~\ref{tab:rlablation}, which has a genuine optimum that moves with scale.

\paragraph{Summary.}
The online formulation changes two conclusions and confirms one. It changes the
status of the Hungarian optimum, which is not an achievable target under churn and
is unboundedly bad if deployed as a policy. It changes what a practitioner should
build: failover is nearly all of the benefit, and the additional complexity of
degradation-aware migration is justified only when degradation persists for longer
than $2M/\mathrm{gap}$ rounds, a quantity \S\ref{sec:online:M} measures for a
real flow table. And it confirms the
paper's central claim in a setting where the claim now has a proof attached: the
sophisticated policy loses to the simple one, and Theorem~\ref{thm:snapshot}
explains why it must.
